% =============================================================
% 4.5 การพัฒนาเกม
% โครงงาน: การออกแบบและพัฒนาเกมเดินเรือด้วยระบบสั่งการด้วยเสียงเพื่อส่งเสริมทักษะการสื่อสาร
% =============================================================

\section{การพัฒนาเกม}

การพัฒนาเกมดำเนินการบนเอนจิน Unity 3D และใช้ภาษา C\# เป็นหลัก โดยแบ่งการพัฒนาออกเป็นระบบย่อยและการบูรณาการเทคโนโลยีปัญญาประดิษฐ์เข้ากับกลไกของเกม ดังนี้

\subsection{การพัฒนาระบบรับคำสั่งเสียง (Speech Recognition System)}
ระบบรับคำสั่งเสียงถูกพัฒนาโดยการเชื่อมต่อไมโครโฟนเข้ากับ API หรือไลบรารีสำหรับการรู้จำเสียงอย่าง Whisper\cite{aguirre2023} เพื่อทำหน้าที่รับสัญญาณเสียงพูดภาษาอังกฤษของผู้เล่นแบบเรียลไทม์ และแปลงข้อมูลเสียงให้เป็นข้อความ (Speech-to-Text) จากนั้นจะส่งข้อความที่ได้เข้าสู่ระบบคัดกรองคำสั่ง (Text Parsing) เพื่อเตรียมนำไปประมวลผลบังคับเรือต่อไป

\begin{figure}[H]
\centering
\includegraphics[width=0.8\textwidth]{Document/GameDev/images/forth.png}
\caption{รหัสคำสั่งส่วนการรับคำสั่งเสียงและการแปลงเป็นข้อความ}
\label{fig:code_speech}
\end{figure}

\subsection{การพัฒนาระบบลูกเรือปัญญาประดิษฐ์ (AI Crew System)}
ระบบนี้ใช้บริการเชื่อมต่อ API ของ OpenAI (เช่น GPT-4o mini) เข้ามาช่วยในการตีความและโต้ตอบ เมื่อระบบได้รับข้อความคำสั่งจากผู้เล่น AI จะทำการตรวจสอบความถูกต้องของบริบทและไวยากรณ์ หากคำสั่งถูกต้อง AI จะส่งคำสั่งไปควบคุมตัวเรือ แต่หากคำสั่งไม่ชัดเจน AI จะทำการสร้างบทสนทนาโต้ตอบกลับ (Voice Feedback) เพื่อขอให้ผู้เล่นย้ำคำสั่ง

\begin{figure}[H]
\centering
\includegraphics[width=0.8\textwidth]{Document/GameDev/images/fifth.png}
\caption{รหัสคำสั่งส่วนการประมวลผลคำสั่งลูกเรือ AI และการโต้ตอบ}
\label{fig:code_ai_crew}
\end{figure}

\subsection{การพัฒนาระบบฟิสิกส์และการควบคุมเรือ (Ship Physics)}
การควบคุมเรือถูกเขียนด้วยภาษา C\# เพื่อจำลองการเคลื่อนที่ของเรือขนาดใหญ่ โดยมีการคำนวณตัวแปรด้านความเฉื่อย (Inertia) และรัศมีวงเลี้ยว (Turning Radius) ทำให้เรือไม่ตอบสนองต่อการเลี้ยวหรือการเบรกในทันทีที่ได้รับคำสั่ง 

\begin{figure}[H]
\centering
\includegraphics[width=0.8\textwidth]{Document/GameDev/images/sixth.png}
\caption{รหัสคำสั่งส่วนระบบฟิสิกส์และการจำลองการเดินเรือ}
\label{fig:code_physics}
\end{figure}

\subsection{ระยะการพัฒนาเกม (Development Phases)}
การดำเนินการพัฒนาเกมตัวต้นแบบถูกแบ่งออกเป็น 4 ระยะหลัก เพื่อให้การทำงานเป็นระบบและสามารถทดสอบได้ในแต่ละขั้นตอน ดังนี้

\begin{enumerate}[label=\arabic*)]
    \item \textbf{ระยะที่ 1 (Prototype):} พัฒนาระบบพื้นฐาน โดยเน้นการเชื่อมต่อระบบไมโครโฟนเข้ากับ Speech-to-Text API และทดสอบการแปลงเสียงเป็นข้อความคำสั่ง (Text Parsing) เพื่อขยับวัตถุเรือเบื้องต้น
    \item \textbf{ระยะที่ 2 (Alpha):} พัฒนาระบบสภาพแวดล้อมทางทะเล การสุ่มเกิดของสิ่งกีดขวาง (Obstacle Generation) และใส่ระบบเสียงโต้ตอบ (Voice-over) ของลูกเรือ AI
    \item \textbf{ระยะที่ 3 (Beta):} พัฒนาระบบ UI/UX ให้สมบูรณ์ ปรับจูนความยากง่าย (เช่น ความเร็วของเรือ ระยะห่างของสิ่งกีดขวาง) และความแม่นยำ (Tolerance) ในการรับคำสั่งเสียง
    \item \textbf{ระยะที่ 4 (Public/Release Candidate):} ตรวจสอบและแก้ไขข้อผิดพลาด (Bug Fixing) และเตรียมแพ็กเกจเกมสำหรับนำไปทดสอบกับผู้ใช้งานจริง
\end{enumerate}